%%%
%
% $Autor: Siddhanth Sharma $
% $Datum: 2026-01-02 $
%
% !TeX spellcheck = GB
% !TeX program = pdflatex
% !TeX encoding = utf8
%
%%%

\chapter{Development Environment -- Hardware / OS / System Tools}
\index{Development Environment -- Hardware}

\section{Version}
\noindent
The project is deployed on the \textbf{Seeed Studio reComputer J1010} platform, integrating an \textbf{NVIDIA Jetson Nano (4~GB, B01 revision)} module on the \textbf{J101 carrier board}.
The OS and NVIDIA acceleration stack are provided via JetPack/L4T.

\vspace{1mm}
\noindent
\textbf{System fingerprint (recorded on-device):}
\begin{itemize}
	\item \textbf{JetPack/L4T:} \texttt{R32.7.1} (JetPack meta package: \texttt{nvidia-jetpack 4.6.2-b5}).
	\item \textbf{Ubuntu:} Ubuntu 18.04.6 LTS (bionic).
	\item \textbf{Kernel:} \texttt{Linux 4.9.253-tegra} (aarch64).
	\item \textbf{CUDA:} \texttt{CUDA 10.2} (\texttt{nvcc} release 10.2, V10.2.300).
	\item \textbf{TensorRT:} \texttt{TensorRT 8.2.1.8} (\texttt{tensorrt==8.2.1.8}).
\end{itemize}

\section{Description}
\noindent
The Jetson system is used as an \textbf{edge AI device} to run the full pipeline locally: camera capture, license plate detection, OCR, and decision/logging.
Operating locally reduces latency and avoids dependence on external servers for the core recognition/billing decision.

\vspace{1mm}
\noindent
Project-relevant system aspects include:
\begin{itemize}
	\item \textbf{Access mode:} local GUI (monitor/keyboard) or headless operation via SSH.
	\item \textbf{Remote development:} VS Code Remote is blocked on JetPack 4.x because the VS Code Remote Server requires \texttt{glibc 2.28}, while Ubuntu 18.04 on Jetson provides \texttt{glibc 2.27}. As a workaround, the project was edited locally and synchronized to the Jetson via SFTP (see \FILE{Code/LicencePlateDetection/.vscode/sftp.json}) while using SSH for terminal commands.
	\item \textbf{Resource constraints:} Jetson Nano has 4\,GB unified memory, requiring memory-aware configuration for stable inference.
\end{itemize}

\section{Installation}
\noindent
The base system installation consists of installing a JetPack/L4T image to the boot medium and completing first-boot initialization.
For this project, the root filesystem is on internal flash storage (\texttt{/dev/mmcblk0p1}, 14\,GB partition), and an external SSD is mounted at \FILE{/media/myssd} for project storage.

\vspace{1mm}
\noindent
A concise installation procedure is:
\begin{itemize}
	\item Flash JetPack/L4T image to the boot medium.
	\item First boot (user creation, locale/time, network).
	\item Update system packages and verify baseline readiness (camera, CUDA/TensorRT where applicable).
\end{itemize}

\section{Configuration}
\noindent
System configuration steps relevant to this project:
\begin{itemize}
	\item \textbf{SSH enablement:} verify headless access for deployment and debugging.
	
	\item \textbf{Performance monitoring:} confirm availability of \SHELL{tegrastats} for runtime observation.
	
	\item \textbf{Power/performance mode:}  
	for benchmarking and stable inference timing, the Jetson is configured to maximum performance using
	\SHELL{sudo nvpmodel -m 0} and \SHELL{sudo jetson\_clocks}.
	
	\item \textbf{Storage:}  
	external SSD mounted at \FILE{/media/myssd} (\texttt{/dev/sda2}, 1.8\,TB).
	
	\item \textbf{Swap:}  
	swapfile at \FILE{/media/myssd/swapfile} (8\,GB) and zram swap devices
	(\texttt{/dev/zram0..3}).
	
	\item \textbf{Camera:}  
	USB camera available as \FILE{/dev/video0} (Microsoft LifeCam).
	
	\item \textbf{TensorRT engine build (on-device):}  
	The deployed Flask dashboard expects a prebuilt TensorRT engine at a fixed location.
	
	{\sloppy\small
		\begin{itemize}
			\item Expected engine path:  
			\FILE{/media/myssd/jetson-inference/data/networks/best.engine}
			
			\item Loader script:  
			\FILE{Code/LicencePlateDetection/jetson-inference/python/www/flask/dashboard2.py}
			
			\item Engine generation:  
			The engine is built from an exported ONNX model using TensorRT’s \SHELL{trtexec} utility:
			
			\SHELL{trtexec --onnx=best.onnx --saveEngine=/media/myssd/jetson-inference/data/networks/best.engine --fp16}
		\end{itemize}
	}
\end{itemize}

\section{First Steps (GUI / UI)}
\noindent
Initial system validation is performed either via GUI terminal or SSH:
\begin{itemize}
	\item Verify OS/L4T identity and network connectivity.
	\item Verify camera detection (presence of \FILE{/dev/video*} and successful test capture).
	\item Verify monitoring tools (\SHELL{tegrastats}) and basic resource availability (RAM/swap/disk).
\end{itemize}

\section{Program ``Hello World'' / Bash Verification}
\noindent
As a minimal OS-level verification, a short Bash script checks whether the platform exposes the required system capabilities
(OS identity, CUDA/TensorRT presence, monitoring tools, and basic resources).

\lstinputlisting[
language=bash,
caption={OS-level ``Hello World'' verification script for the Jetson platform},
label={lst:jetson_hello_world_bash},
basicstyle=\ttfamily\small,
keywordstyle=\color{blue},
commentstyle=\color{gray},
stringstyle=\color{green!40!black},
breaklines=true,
showstringspaces=false,
numbers=left,
numberstyle=\tiny\color{gray},
frame=single,
rulecolor=\color{black!30}
]{../../Code/SystemTools/jetson_hello_world.sh}